The register lists the 191 computable, nonconstant indicator items used in the association analysis in Section 4.2 of the main paper. Items are grouped into the same five categories as the method. An item is retained when at least one numerical form has two distinct finite values and positive scenario-weighted variation; individual forms may be constant in other decision types. Inclusion does not imply statistical significance or selection by every prediction model. Section~\ref{supp:other-indicators} defines the remaining catalog items and records why no association is available.

Arrival and departure descriptors, train-pair contrasts, and summaries over different time windows provide the numerical forms. Here $W$ denotes a time interval centered on a reference departure, with a total width of 5, 10, or 20 minutes. Blockage start and reopening use $\tau_b$ and $\tau_e$, as in the formulation. Summary measures include counts, totals, means, extrema, dispersion, and quantiles. Missing measurements remain distinct from zero.

Definitions reflect the information available in the reference timetable. H17--H19 measure competition within a time window, while P15 uses track count as an overtaking-opportunity proxy. Identifiers are used only to match records and resolve tied prediction scores.

\begingroup
\small
\setlength{\tabcolsep}{3pt}
\renewcommand{\arraystretch}{1.06}
\begin{longtable}{@{}>{\raggedright\arraybackslash}p{0.8cm}>{\raggedright\arraybackslash}p{3.7cm}>{\raggedright\arraybackslash}p{10.0cm}@{}}
\caption{Computable, nonconstant indicator catalog}\label{tab:indicator-register}\\
\toprule Code & Indicator & Operational definition \\ \midrule
\endfirsthead
\multicolumn{3}{l}{\textit{Table \thetable\ continued}}\\
\toprule Code & Indicator & Operational definition \\ \midrule
\endhead
\midrule\multicolumn{3}{r}{\textit{Continued on next page}}\\
\endfoot
\bottomrule
\multicolumn{3}{@{}p{14.5cm}@{}}{\footnotesize\textit{Notes:} Codes identify catalog items. Only computable items with nonconstant numerical forms in the association analysis of Section 4.2 of the main paper are listed. The symbols $\tau_b$ and $\tau_e$ denote blockage start and reopening. The interval $W$ is centered on a reference departure, with a total width of 5, 10, or 20 minutes as specified by the numerical form.}\\
\endlastfoot
\multicolumn{3}{@{}l}{\textbf{Disruption conditions}}\\*[2pt]
D01 & Blockage start time & Blockage start $\tau_b$, in minutes after 08:00. \\
D02 & Blockage end & Known reopening time $\tau_e$, in minutes after 08:00. \\
D03 & Blockage duration & $\tau_e$ minus $\tau_b$, in minutes. \\
D04 & Start-time wave phase & $\tau_b$ modulo the 40-minute departure-wave period. \\
D05 & Reopening wave phase & $\tau_e$ modulo the 40-minute departure-wave period. \\
D08 & Remaining blockage & Known remaining blockage duration at the initial decision. \\
D09 & Time from start time & Reference or planned departure time minus $\tau_b$. \\
D10 & Time from reopening & Reference or planned departure time minus $\tau_e$. \\
D11 & Blocked section & Separate indicators for H1--C, U13--H1 and U11--U12; categorical location, not a continuous distance. \\
D20 & Lost resource time & Two blocked directional running resources times blockage duration. \\
E01 & Train state at start time & State indicators derived from the original timetable and committed history. \\
E02 & Executed route fraction & Fraction of station departures already executed by start time. \\
E03 & Current route position & Route index of the emergency stopping position or completed route end. \\
E04 & Committed running remainder & Remaining time of an irrevocable run to the emergency station. \\
E05 & Emergency stopping position & Route index of the first station reached under the emergency stopping rule. \\
E06 & Time to emergency station & Known time from start time to the emergency station. \\
E07 & Emergency arrival time & Known emergency-station arrival time, using $\tau_b$ for an origin hold. \\
E08 & Minimum emergency hold & Required hold from emergency arrival or start time until $\tau_e$. \\
E09 & Already stopped at start time & Indicator that the train is at a station when blockage begins. \\
E12 & Origin waiting count & Network-wide number of emergency origin holds. \\
E13 & Station waiting count & Network-wide number of emergency holds away from origins. \\
E16 & Emergency backlog & Number of trains subject to the emergency stopping rule. \\
E17 & Early release concentration & Fraction of emergency trains released before $\tau_e$ plus five minutes. \\
E19 & Committed event count & Number of arrival and departure events fixed by operational history. \\
E20 & Planned starts before release & Number of trains planned to start strictly between $\tau_b$ and $\tau_e$. \\
E21 & Train-state counts at blockage start & Counts of origin holds, committed-run completions, station holds and completed trains. \\
E26 & Committed movement count & Recorded runs that must finish despite the emergency rule. \\
I05 & Predicted remaining duration & Known remaining blockage duration under the deterministic assumption. \\
I06 & Remaining-duration bounds & Lower and upper predicted remaining durations; both equal the known duration here. \\
I09 & Remaining-duration quantiles & The 0.1, 0.5 and 0.9 quantiles of remaining duration; equal under deterministic recovery. \\
X01 & Blocked-section exposure & Whether the train's fixed route includes the blocked section. \\
X02 & Already past blocked section & Whether the planned exit from the blocked section is no later than $\tau_b$. \\
X03 & Planned clearance before start time & Whether the train was planned to clear the blocked section by $\tau_b$. \\
X04 & Planned blockage overlap & Length of the planned blocked-section occupation intersecting $[\tau_b,\tau_e)$. \\
X05 & Planned overlap fraction & Blockage overlap divided by planned blocked-section running duration. \\
X06 & Planned entry from start time & Planned entry to the blocked section minus $\tau_b$. \\
X07 & Planned exit from start time & Planned exit from the blocked section minus $\tau_b$. \\
X08 & Sections to blockage entry & Remaining section count from the emergency position to blockage entry. \\
X09 & Minimum time to entry & Sum of minimum running times from the emergency position to blockage entry. \\
X10 & Signed blockage distance & Blockage-entry route index minus the candidate station index. \\
X11 & At blockage entrance & Whether the candidate station is the blocked-section entrance. \\
X12 & Can clear before start time & Single-train minimum-running estimate of clearing the blocked section before $\tau_b$. \\
X13 & Arrival before reopening & Whether the emergency waiting station can be reached before reopening under the recorded operating state. \\
X14 & Earliest-entry estimate & Single-train entry-time estimate, ignoring other trains. \\
X15 & Earliest entry from release & Earliest-entry estimate minus $\tau_e$. \\
X16 & Blockage-only waiting & Nonnegative $\tau_e$ minus earliest-entry estimate for a train not cleared by start time. \\
X17 & Remaining route fraction & Remaining section count divided by total route section count. \\
X18 & Emergency origin hold & Whether the emergency stopping rule holds this train at its origin. \\
X19 & Emergency restriction type & Recorded category: no emergency hold, origin hold, committed-run completion, or station hold. \\
\multicolumn{3}{@{}l}{\textbf{Timetable characteristics}}\\*[2pt]
F12 & Minimum dwell & Station-specific minimum dwell in the operational input. \\
F13 & Minimum running time & Train- and section-specific minimum running duration. \\
F14 & Allowed stopping mode & Required-stop and unfixed optional-stop indicators in the operational input. \\
M03 & Original domain width & Raw operational upper minus lower bound before common propagation. \\
M04 & Propagated domain & Lower bound, upper bound, domain width or order-choice count after common preprocessing. \\
M05 & Reference adjustment room & Reference value minus lower bound, upper bound minus reference, or reference order. \\
M06 & Remaining candidate decisions & Counts of the train's remaining stopping, duration, precedence and track candidates. \\
T01 & Planned arrival time & Original arrival time at the station. \\
T02 & Planned departure time & Original visit departure time. \\
T03 & Planned dwell & Original departure minus original arrival at a station. \\
T04 & Planned running duration & Original duration of the outgoing running activity. \\
T05 & Planned stopping mode & Original planned stopping indicator, without treating required stops as cancellable. \\
T06 & Planned track label & Recorded numeric track-label suffix; an assignment descriptor, not a cardinal capacity. \\
T09 & Planned entry buffer & Gap to the preceding planned entry minus the minimum entry headway. \\
T10 & Planned exit buffer & Gap to the preceding planned exit minus the minimum exit headway. \\
T11 & Planned track buffer & Planned arrival minus the latest preceding release on the same planned track. \\
T12 & Planned local flow & Planned running-resource entries in $W$, divided by window width. \\
T13 & Planned train-type pattern & Fast-train share and consecutive type-alternation fraction in $W$. \\
T14 & Departure wave & Integer origin departure time divided by the 40-minute wave period. \\
T15 & Fast-train indicator & Recorded fast/slow class, distinguished by minimum running duration. \\
T17 & Total planned journey time & Planned terminal arrival minus origin departure. \\
T18 & Remaining planned journey & Planned terminal arrival minus the candidate planned departure. \\
T19 & Remaining minimum journey & Remaining minimum runs plus subsequent mandatory minimum dwells. \\
T20 & Remaining scheduled stops & Count of subsequent stops required by the original timetable. \\
T21 & Nearby scheduled-stop distance & Route-index distance to the preceding and following stops in the original timetable. \\
T22 & Route extent & Total and remaining route section counts. \\
T24 & Planned terminal position & Planned terminal arrival and its difference from $\tau_e$. \\
\multicolumn{3}{@{}l}{\textbf{Train interactions}}\\*[2pt]
B17 & Relation-layer overlap & Jaccard overlap of train-pair sets in entry, exit and station-track layers. \\
E11 & Co-located emergency trains & Number of same-direction emergency trains waiting at the same station. \\
E14 & Reference release rank & Rank in the emergency-station queue ordered by reference release time. \\
E15 & Queue fast-train share & Fast-train fraction in the emergency-station queue. \\
H17 & Window competition size & Distinct trains occupying the station in $W$; a window-based proxy. \\
H18 & Co-window occupations & Number of train occupations at the station intersecting $W$; a competition-membership proxy. \\
H20 & Support distance & Nearest optional-stop route-index distance and finite active distance for an order pair. \\
H21 & Shared-visit coupling & Number of candidate decision classes and order candidates incident on the visit. \\
H22 & Shared downstream support & Jaccard overlap of downstream event sets for endpoints or train pairs. \\
P02 & Planned rank distance & Following-minus-leading rank on the common running resource. \\
P03 & Planned entry gap & Following-minus-leading planned resource-entry time. \\
P04 & Planned exit gap & Following-minus-leading planned resource-exit time. \\
P05 & Running-duration contrast & Following-minus-leading planned running duration on the common resource. \\
P06 & Train-type pair & Fast/fast, fast/slow, slow/fast and slow/slow category indicators. \\
P07 & Remaining-journey contrast & Following-minus-leading remaining planned and minimum journey times. \\
P08 & Terminal-time contrast & Following-minus-leading planned terminal arrival. \\
P09 & Stop-pattern disagreement & Number of shared stations at which planned stopping indicators differ. \\
P10 & Shared route extent & Common section count and the leading train's summed minimum time on those sections. \\
P11 & Shared station pools & Number of stations common to both routes. \\
P12 & Merge and split positions & First and last shared-station indices relative to the candidate visit. \\
P13 & Catch-up proxy & Positive minimum-running-time advantage over common sections minus planned entry gap. \\
P14 & Potential order positions & Number of free precedence candidates involving this train pair. \\
P15 & Parallel common stations & Number of common stations with at least two same-direction tracks. \\
P16 & Reference-delay contrast & Following-minus-leading reference departure delay. \\
P17 & Same reference track & Whether both trains use the same nonempty track label in the reference. \\
P18 & Reference interval overlap & Overlap of the two reference running intervals on the common resource. \\
X20 & Same exposure category & Whether the two precedence trains have the same emergency restriction category. \\
X21 & Pair entry-time contrast & Following-minus-leading earliest-entry time relative to $\tau_e$. \\
X22 & Pair waiting contrast & Following-minus-leading blockage-only waiting estimate. \\
X23 & Pair blockage-rank gap & Planned rank difference of the two trains on their common running resource. \\
\multicolumn{3}{@{}l}{\textbf{Infrastructure utilization}}\\*[2pt]
B07 & Regional saturation & Station saturation fraction in $W$. \\
B13 & Adjacent-window difference & Following minus preceding window mean occupation. \\
E10 & Emergency pool capacity & Number of same-direction tracks at the assigned emergency station. \\
E18 & Parking concentration & Herfindahl concentration of emergency holds across stations. \\
E23 & Initial station occupation & Same-direction station occupation at $\tau_b$ in the recorded reference. \\
F01 & Facility role & Hub and common-corridor indicators for the candidate station. \\
F02 & Physical degree & Degree in the undirected physical station graph. \\
F03 & Branch-count proxy & Physical station degree minus one, bounded below by zero. \\
F04 & Physical betweenness & Unweighted betweenness in the physical station graph. \\
F05 & Physical articulation point & Whether the station is an articulation point of that graph. \\
F06 & Directional pool size & Number of same-direction station tracks. \\
F07 & Compatible tracks & Number of tracks allowed for the visit; zero at external boundaries. \\
F15 & Parallel-capacity proxy & Same-direction track-pool size as a proxy for parallel reception. \\
F16 & Downstream capacity gradient & Next station's directional capacity minus current capacity. \\
F17 & Distance to a hub & Minimum physical section count to H1 or H2. \\
F18 & Next stopping opportunity & Route-section distance to the next station at which stopping is permitted. \\
H19 & Peak demand/capacity & Peak same-direction occupation in $W$ divided by track-pool size. \\
R01 & Mean occupation & Time-average number of occupied same-direction station tracks in $W$. \\
R02 & Peak occupation & Maximum same-direction station occupation in $W$. \\
R03 & Minimum free capacity & Directional capacity minus peak occupation in $W$. \\
R04 & Utilization & Mean occupation divided by directional track-pool capacity. \\
R05 & Saturated fraction & Fraction of $W$ at which occupation reaches directional capacity. \\
R06 & Saturation persistence & Longest consecutive saturated interval within $W$. \\
R07 & Entry load & Reference entries in $W$ times minimum entry headway, divided by width. \\
R08 & Exit load & Reference exits in $W$ times minimum exit headway, divided by width. \\
R09 & Reference headway buffer & Reference entry/exit gap to the preceding train minus minimum headway. \\
R10 & Intersecting dwell distribution & Durations of train occupations at the station whose occupations intersect $W$; not passenger waiting. \\
R11 & Flow imbalance & Station arrivals minus departures within $W$. \\
R12 & Demand/capacity proxy & Mean station occupation divided by directional capacity. \\
R13 & Adjacent-station occupation & Upstream/downstream occupation at reference departure and their difference. \\
R14 & Window competitors & Distinct trains with station occupations intersecting $W$. \\
R15 & Free windows & Count, total duration and maximum duration of unsaturated intervals in $W$. \\
R16 & Planned-arrival overlap & Planned station arrivals in $W$ that overlap a reference occupation. \\
R17 & Occupation change rate & Difference of following and preceding window means divided by twice the width. \\
R18 & Opposite-direction utilization & Opposite-direction mean occupation normalized by the corresponding capacity here. \\
\multicolumn{3}{@{}l}{\textbf{Delay propagation patterns}}\\*[2pt]
B01 & Local region size & Same-station, same-direction event and distinct-train counts in $W$. \\
B02 & Local edges and density & Internal edge count and directed density in the local region. \\
B03 & Active local density & Active internal edge count divided by possible directed event pairs. \\
B04 & Boundary connections & Numbers of edges entering and leaving the local region. \\
B05 & Regional delay distribution & Reference event-delay summaries for the local region. \\
B06 & Regional delay variance & Variance of reference event delays in the local region. \\
B08 & Community structure & Greedy-modularity community size and links to other communities. \\
B09 & Boundary pressure imbalance & Outgoing minus incoming delay support across the region boundary. \\
B10 & Affected-train share & Fraction of trains with positive reference terminal delay. \\
B11 & Global path distribution & Summaries of longest active path lengths. \\
B12 & Delay concentration & Entropy and Herfindahl indices of delay totals by station and 20-minute time bin. \\
B15 & Global degree distribution & Incoming/outgoing active-degree summaries over graph events. \\
B16 & Global reachability & Reachable ordered-event-pair fraction and mean finite active distance. \\
G01 & Reference arrival delay & Reference arrival minus original planned arrival. \\
G02 & Reference departure delay & Reference departure minus original planned departure. \\
G03 & Activity delay increment & Difference of reference endpoint delays for dwell or running activities. \\
G04 & Activity-duration increment & Reference dwell/running duration minus its planned duration. \\
G05 & Event degree & Incoming and outgoing edge counts in the reference event graph. \\
G06 & Cross-train degree & Incoming and outgoing edges linking different trains. \\
G07 & Degree by relation & Incoming/outgoing counts for dwell, run, entry, exit, track and release relations. \\
G08 & Active degree & Numbers of active incoming and outgoing edges. \\
G09 & Active cross-train links & Count and fraction of active links between different trains. \\
G10 & Edge-slack distribution & Incoming/outgoing edge slack summaries. \\
G11 & Near-tight links & Incident edges with slack at most one or two minutes. \\
G12 & Delay-support distribution & Incoming/outgoing positive delay-support summaries. \\
G13 & Upstream-neighbor delay & Delay summaries of directly upstream different trains, deduplicated by train. \\
G14 & Downstream-neighbor delay & Delay summaries of directly downstream different trains, deduplicated by train. \\
G15 & Active-neighbor dispersion & Range, variance and coefficient of variation of active-neighbor delays. \\
G16 & Incident-link imbalance & Outgoing minus incoming degree and delay-support sums. \\
G17 & Relation-type entropy & Entropy of incident edge-type counts. \\
G18 & External-source share & Fraction of incoming edges from the merged external release source. \\
G19 & Endpoint contrasts & Departure/arrival feature differences and recorded running-endpoint contrasts. \\
H01 & Shortest source distance & Shortest active-source edge count and recorded cumulative minimum-lag distance. \\
H02 & Longest active path & Longest active directed path ending at the event. \\
H03 & Event time from release & Reference event time minus reopening time $\tau_e$. \\
H04 & Upstream reachability & Number of active upstream events, excluding the merged source. \\
H05 & Downstream reachability & Number of active downstream events. \\
H06 & Upstream trains & Distinct trains among active upstream events. \\
H07 & Downstream trains & Distinct trains among active downstream events. \\
H08 & Spatial reachability & Distinct stations and branches reachable upstream/downstream. \\
H09 & Reachable delay totals & Event-delay sums and sums of per-train maximum delays upstream/downstream. \\
H10 & Active betweenness & Normalized unweighted betweenness in the active directed graph. \\
H11 & Harmonic closeness & Normalized sums of reciprocal finite active distances, incoming and outgoing. \\
H12 & PageRank & Active-graph PageRank with damping 0.85. \\
H13 & Weak-component size & Event and distinct-train counts in the active weak component. \\
H15 & Chain and crossing length & Longest active chain and maximum number of cross-train edges along a chain. \\
H16 & Local motifs & Fan-in pairs, fan-out pairs and directed diamonds at their source. \\
H23 & Local clustering & Local clustering coefficient in the undirected projection of the active propagation graph. \\
H24 & Projection core number & Core number in the active graph's undirected projection. \\
H25 & Local assortativity & Neighbor same-speed-class fraction and mean absolute delay difference. \\
\end{longtable}
\endgroup


\subsection{Retrieved studies used to develop the catalog}
The 22 studies below supplied the literature context for developing the candidate indicator catalog. This list records the consulted studies; it does not imply that every candidate indicator was proposed in a cited study or that full texts were available for every study. Study identifiers and the original metadata are also provided in the accompanying retrieved-studies CSV and JSON files.
\begingroup
\small
\setlength{\tabcolsep}{5pt}
\begin{longtable}{@{}p{.04\linewidth}>{\raggedright\arraybackslash}p{.20\linewidth}>{\raggedright\arraybackslash}p{.68\linewidth}@{}}
\caption{Retrieved studies used to develop the indicator catalog}\label{tab:retrieved-studies}\\
\toprule No. & Study & Title and source \\ \midrule
\endfirsthead
\toprule No. & Study & Title and source \\ \midrule
\endhead
\bottomrule
\endfoot
1 & Goverde (2007) & Railway timetable stability analysis using max-plus system theory. \href{https://doi.org/10.1016/j.trb.2006.02.003}{Source} \\
2 & Goverde (2010) & A delay propagation algorithm for large-scale railway traffic networks. \href{https://doi.org/10.1016/j.trc.2010.01.002}{Source} \\
3 & Büker and Seybold (2012) & Stochastic modelling of delay propagation in large networks. \href{https://doi.org/10.1016/j.jrtpm.2012.10.001}{Source} \\
4 & Harrod, Cerreto and Nielsen (2019) & A closed form railway line delay propagation model. \href{https://doi.org/10.1016/j.trc.2019.02.022}{Source} \\
5 & Andersson, Peterson and Törnquist Krasemann (2013; online 2014) & Quantifying railway timetable robustness in critical points. \href{https://doi.org/10.1016/j.jrtpm.2013.12.002}{Source} \\
6 & Ma et al. (2021) & Scheduling-measure dependent modelling of delay propagation on a single railway line: A matrix transformation approach. \href{https://doi.org/10.1049/itr2.12112}{Source} \\
7 & Meng et al. (2019) & Increasing Robustness by Reallocating the Margins in the Timetable. \href{https://doi.org/10.1155/2019/1382394}{Source} \\
8 & Wu et al. (2023) & Cascading Delays for the High-Speed Rail Network Under Different Emergencies: A Double Layer Network Approach. \href{https://doi.org/10.1109/JAS.2022.105530}{Source} \\
9 & Lian et al. (2024) & Bi-Directional Delay Propagation Analysis and Modeling for High-Speed Railway Networks Under Disturbance. \href{https://doi.org/10.1109/TCSS.2023.3322747}{Source} \\
10 & Tessitore et al. (2025) & On the fragility of a train timetable. \href{https://doi.org/10.1016/j.omega.2025.103341}{Source} \\
11 & Corman et al. (2014) & Evaluating Disturbance Robustness of Railway Schedules. \href{https://doi.org/10.1080/15472450.2013.801714}{Source} \\
12 & Zhang et al. (2024) & Propagation-Based Train Rescheduling under Recoverable Delay Disturbances. \href{https://doi.org/10.1155/2024/9961840}{Source} \\
13 & Huang et al. (2026; online 2025) & Solving the railway timetable rescheduling problem with graph neural networks. \href{https://doi.org/10.1007/s40534-025-00383-7}{Source} \\
14 & Acuña-Agost et al. (2011) & SAPI: Statistical Analysis of Propagation of Incidents. A new approach for rescheduling trains after disruptions. \href{https://doi.org/10.1016/j.ejor.2011.05.047}{Source} \\
15 & Tessitore et al. (2025) & An iterative fragility-based approach to timetable design. \href{https://doi.org/10.1016/j.jrtpm.2025.100554}{Source} \\
16 & Hua et al. (2024) & Improving timetable resilience for high-speed railways through time supplement allocation. \href{https://doi.org/10.11860/j.issn.1673-0291.20230073}{Source} \\
17 & Corman and Kecman (2018) & Stochastic prediction of train delays in real-time using Bayesian networks. \href{https://doi.org/10.1016/j.trc.2018.08.003}{Source} \\
18 & Ulak et al. (2020) & Analyzing network-wide patterns of rail transit delays using Bayesian network learning. \href{https://doi.org/10.1016/j.trc.2020.102749}{Source} \\
19 & Dekker et al. (2022) & Modelling railway delay propagation as diffusion-like spreading. \href{https://doi.org/10.1140/epjds/s13688-022-00359-1}{Source} \\
20 & Ge et al. (2024; online 2023) & Modelling the cascading effects of train delay patterns and inter-train control actions with Bayesian networks. \href{https://doi.org/10.1080/23248378.2023.2194304}{Source} \\
21 & Daniotti et al. (2024) & Systemic risk approach to mitigate delay cascading in railway networks. \href{https://doi.org/10.1038/s44333-024-00012-6}{Source} \\
22 & Xu et al. (2026) & Multivariate event hypergraph diffusion model for train delay prediction. \href{https://doi.org/10.1016/j.trc.2025.105390}{Source} \\
\end{longtable}
\endgroup

